\section{Introduction}

\subsection{Main results}

Let \((\Gamma,\cS)\) be a countable group with a symmetric finite set of generators \(S\).

For \(d \ge 2\), consider \(P\) a subset of \(\lbrace 1,\ldots, d-1\rbrace\).
A representation \(\rho: \Gamma \to \SL(d,\R)\)
is called a {\it{\(P\)-Anosov}} representation if there exists \(c > 0\) such that
\[
\frac{s_{p}(\rho(\gamma))}{s_{p+1}(\rho(\gamma))} > c\exp(c|\gamma|),
\]
for all \(\gamma \in \Gamma\) and \(p \in P\). Here \(|\gamma|\) denotes word length in \(\Gamma\) with respect to \(S\) and, for~\(1\le i\le d, \, s_i(\g)\) denotes the \(i\)-th largest singular value of \(\g \in \SL(d,\R)\) with respect to the usual inner product on \(\R^d\). A representation is called {\it{Borel-Anosov}} if it is \(\lbrace 1,\ldots, d-1\rbrace\)-Anosov.

Anosov representations have been introduced by F.\,Labourie (\cite{Lab06}) and have been the subject of many studies (\cf in particular \cite{GGKW17}, \cite {KLP16}, \cite{KLP17}, \cite{KLP18a}, \cite{BPS19} and \cite{canary} for a recent survey). It is widely accepted that non-trivial Anosov representations present the right analog of convex cocompact representations in higher rank. Many properties are suggested by this analogy. In another direction, we~present in this note a phenomenon which is purely higher rank: there is a dimension gap for minimal invariant subsets for the action of \(\rho (\G) \) on the spaces of flags. The proof uses a new variational principle for the action of \(\rho (\G) \) on the spaces of partial flags (see below Theorem \ref{theoremb}).

Let \(Q\) be a non-empty subset of \(\lbrace 1,\ldots, d-1 \rbrace\). We~consider the space \(\F_Q\) of partial flags with signature \(Q\), endowed with a rotationally invariant Riemannian metric. Let~us recall the definition of the Minkowski dimension of a set \(\La \subset \F_Q\): for \(\e >0\), let~\(N(\La, \e) \) be the covering number of the set \(\La \) by balls of radius \(\e\) in the metric of~\(\F_Q\). The {\it {Minkowski (or Box) dimension }} \(\dim_B(\La) \) is defined by
\[
 \dim_B (\La) \; : =\; \limsup_{\e \to 0} \frac{\log N(\La, \e)}{\log (1/\e)}.
\]
Recall that the Hausdorff dimension of a metric set is bounded above by its Box dimension. Our result is
\begin{theorem}[Dimension gap for minimal sets of Anosov representations]\label{main}
With the above notations, let \(\La\) be a minimal invariant set for the action of \(\rho (\G) \) on \(\F_Q\). Assume that the representation \(\rho\) is \(P\)-Anosov for some \(P \subset \lbrace 1,\ldots, d-1 \rbrace\) and such that \(\rho (\G)\) is Zariski dense in \(\SL(d,\R)\). Then, with \(M := \# (Q \cap P)\),
\[
\dim_B(\Lambda) \le \dim(\F_Q) - \frac{M - 1}{2}.
\]
\end{theorem}

In the case \(d = 2\) the bound given by the above theorem is trivial. And, indeed, uniform discrete subgroups of \(\SL(2,\R)\) have full limit set. A dimension gap for the limit set of classical Schottky groups was proved in \cite{Doy88}. We~note, however, that there exist uniform discrete (and thus convex co-compact) subgroups of isometries of
rank one symmetric spaces and thus there is no dimension gap for convex cocompact representations in rank one. On the other hand, there are many conditions that ensure that the limit set is a Lipschitz submanifold of the flag space (see the discussions in \cite{PSW19}). Our result is easy in those cases.

If \(\rho\) is \(P\)-Anosov then it is also \(P \cup P^\ast\)-Anosov for \(P^\ast = \lbrace d-p: p \in P\rbrace\). Theorem~\ref{main} gives a non-trivial upper bound for the dimension of any minimal invariant set in \(\F_{Q}\) for all \(Q\) with \(\#(Q\cap (P \cup P^\ast)) \ge 2\). In particular, as soon as the representation is \(k\)-Anosov for some \(k \not = d/2\) and Zariski dense, then the codimension of the limit set in the full flag space is at least 1/2.

There are a few simple reductions which allow us to state our core result. Firstly, we~observe that the bundle \(\pi_{Q, Q\cap P} : \F_Q \to \F_{Q\cap P} \) is smooth, so that for any closed \(\La \subset \F_{Q\cap P}, \) we have
\[
 \dim_B ((\pi_{Q, Q\cap P})^{-1} (\La)) \; = \; \dim_B (\La) + \dim \F_Q - \dim \F_{Q\cap P}.
\]
So, it suffices to show Theorem \ref{main} for the case \(Q= P\). In that case, the unique minimal invariant subset of \(\F_{P}\) is the limit set \(\La_\rho\).

By considering \(\G /{\textrm {Ker}} \rho\), we may assume that the representation \(\rho\) is faithful. Furthermore, we~may assume that \(\Gamma\) is non-elementary hyperbolic since, on the one hand, an Anosov subgroup of \(\SL (d,\R) \) is hyperbolic (\cite{KLP18}, \cite[Th.\,3.2]{BPS19}) and, on~the other hand, the limit set of an Anosov representation of an elementary (\ie virtually cyclic) hyperbolic group is finite. Taking all this into account, we~only have to prove
\begin{theorem}[Dimension gap for limit sets of Anosov representations]\label{maintheorem}
Let \(\G\) be a non-elementary, hyperbolic, finitely generated group, \(P \subset \lbrace 1,\ldots, d -1\rbrace\) with \(M:= \#P \ge 2 \), \(\rho \) a faithful \(P\)-Anosov representation of \(\G\) in \(\SL (d,\R) \) with \(\rho(\G)\) Zariski dense in \(\SL (d,\R) \) and \(\La_\rho \subset \F_{P} \) the limit set of the representation \(\rho\). Then,
\[
\dim_B(\Lambda_\rho) \le \dim(\F_{P}) - \frac{M - 1}{2}.
\]
\end{theorem}

It seems natural to conjecture that a dimension gap will exist between the limit set and the corresponding space of flags, for any Zariski dense representation into a reductive algebraic Lie group of rank \(r \ge 2\) which is Anosov with respect to at least two simple roots.

Let \(g\in \SL(d,\R)\). The numbers \(\la_i (g) := \lim_n \frac{1}{n }\log s_i(g^n) \) are the logarithms of the moduli of the eigenvalues of \(g\). Apply the Oseledets theorem in the trivial case of one matrix
\(g \in \SL (d,\R)\). The Oseledets decomposition is obtained from the Jordan decomposition by grouping together the spaces corresponding to the eigenvalues with the same modulus. These moduli are the \(\exp \la_i (g)\).

In our setting, if~\(\g \in \G\) (but a finite number), for all \(p \in P, \; \la_p (\rho(\g)) -\la_{p+1}(\rho(\g)) \geq c |\g| \). We~associate to every such \(\g \in \G \) an element \(\eta (\rho (\g))\) of the space \(X_P\) of decompositions of \(\R^d = E_1 \oplus \ldots \oplus E_{M+1} \), with \(\dim E_k = p_k- p_{k-1}\), namely
\[
 \eta (\rho (\g)) \; = \; E_1(\rho (\g)) \oplus \ldots \oplus E_{M+1} (\rho (\g)), 
\]
with the property that for all \(k = 1, \ldots, M+1\), the space \(E_{k} (\rho (\g)) \) is generated by the Jordan spaces of \(\rho(\g)\) such that the modulus \(\exp \la\) of the eigenvalue satisfies\footnote{With the convention that \(p_0=0, \la_{p_0} =+\infty\) and \(p_{M+1}:= d\).}
\[
\exp \la_{p_{k-1}}(\rho (\g)) >\exp \la \ge \exp \la_{p_{k}}(\rho (\g)).
\]
Observe that the space \(X_P\) can also be seen as the open \(\SL (d,\R) \)-invariant set of pair of flag spaces in general position in \(\F_P \x \F_{P^\ast} \).

\begin{corollary}\label{eigenvectors} Let \(\G\) be a hyperbolic finitely generated group, \hbox{\(P \!\subset\! \lbrace 1,\ldots, d-1 \rbrace\)} with \(M:= \# P \).
With the above notations, if~\(\rho \) is a \(P\)-Anosov, Zariski dense representation of \(\G\) in \(\SL(d,\R) \), set \(\Om_\rho \subset X_P \) the closure of the set \(\{ \eta (\rho (\g)); \g \in \G\}\). Then,
\[
\dim_B(\Om_\rho) \le \dim X_P -M +1.
\]
\end{corollary}

Indeed, Corollary \ref{eigenvectors} is trivial for \(d= 2\) and for \(M=1\). It is trivial as well if \(\G\) is elementary. For \(d \ge 3\) and \(M\ge 2\), we show in Section \ref{splitting} why Corollary \ref{eigenvectors} follows from Theorem \ref{maintheorem}.

The first non-trivial case for our results is when \(d=3\) and \(P= \{1,2\}\). So,~let~\(\rho \) be a Borel-Anosov representation of the finitely generated non-elementary group \((\G, S)\) in \(\SL(3,\R) \) and we can consider the limit set \(\Lambda_\rho \) of the action of \(\rho(\G) \) on the space~\(\F\) of complete flags in \(\R^3\). Besides, for all \(\g \in \G\), the matrix \(\rho(\g)\) admits three distinct eigenspaces written \((E_1 (\rho(\g)), E_2 (\rho(\g)), E_3 (\rho(\g)))\) in the order of the absolute values of the eigenvalues. Let \(\Om_\rho \) denote the closure in \(\PP^2\x\PP^2\x\PP^2\) of the \(\{(E_1 (\rho(\g)), E_2 (\rho(\g)), E_3 (\rho(\g))), \g \in \G \} \). We~obtain, from Theorem \ref{maintheorem} and Corollary~\ref{eigenvectors}:

\skpt
\begin{corollary}\label{2.5, 5}\mbox{}\vspace*{-\baselineskip}
\[
\dim_B(\Lambda_\rho) \le \frac{5}{2} < 3 = \dim(\F), \; \dim_B(\Om_\rho) \le 5 < 6 = \dim (\PP^2\x\PP^2\x\PP^2).
\]
\end{corollary}
Let \(\G\) be a surface group. For the Hitchin component, the dimension of the projective limit set is 1 (see \cite{Lab06}) and our result is not new. It is is new and unexpected, as far as the authors are aware, for the Zariski dense representations that lie in the same connected component as the Teichmüller times Identity representations, \ie in the Barbot component (\cf \cite{barbot}).

\Subsection{Strategy of the proof of Theorem \ref{maintheorem}}

\subsubsection{Random walk entropy}

Let \(\M\) be the set of probability measures \(\mu\) on \(\Gamma\) with countable support and with finite first moment, meaning
\[
\sum_{\gamma \in \Gamma}\mu(\gamma) |\gamma| < +\infty,
\]
such that the semi-group \(\Gamma_\mu\) generated by the support of \(\mu\) is non-elementary (\ie contains two independent loxodromic elements).
By \cite{KV83}, if~\(\mu \in \M\) then \(\mu\) has finite Shannon entropy, meaning
\[
H(\mu) = -\sum_{\gamma \in \Gamma}\mu(\gamma)\log(\mu(\gamma)) < +\infty.
\]
The random walk entropy \(h_\mu\) of \(\mu \in \M\) is defined by
\begin{equation}\label{rwentropy} h_\mu = \lim_{n \to +\infty}\frac{1}{n}H(\mu^{\ast n}),\end{equation}
where \(\mu^{\ast n}\) is the \(n\)-fold convolution of \(\mu\) with itself.

\subsubsection{Lyapunov exponents}\label{lyapunovexp}
Denote by \(\aaa^+ \) the cone
\[
\ts \aaa^+ := \lbrace a \in \R^{d}: a_1 \ge \dots \ge a_d, \sum a_i = 0\rbrace.
\]
For \(g \in \SL(d,\R)\) one has \(\log S(g) \in \aaa^+\), where
\[
\log S(g): = (\log s_1(g),\ldots, \log s_d(g)). 
\]
The Lyapunov exponents \(\lambda(\rho_*\mu) = \{\lambda_i(\rho_\ast\mu), 1 \le i \le d \} \in \aaa^+\) of the random walk \((\G, \mu)\) in the representation \(\rho\) are defined by
\begin{equation}\label{exponents}
\lambda(\rho_*\mu):= \lim_{n \to +\infty} \frac{1}{n}\sum_{\gamma \in \Gamma}\mu^{\ast n}(\gamma)\log\bigl(S(\rho(\gamma))\bigr).\end{equation}
The limits exist by the sub-additivity of the sequences
\[
\ts \{\sum_{j\le i}\sum_{\gamma \in \Gamma}\mu^{\ast n}(\gamma)\log\bigl(s_j(\rho(\gamma))\bigr) \}_{n \in \N}.
\]

\subsubsection{Lyapunov and Falconer dimension}\label{Lyadimension}

We say a pair \((i,j)\), \(0\!<\!i\!<\!j\!\le\! d\), is~separated by \(p \in P\) if \(i \le p < j\). For each \(p \in P\) let \(S_p\) be the set of pairs separated by \(p\), and define \(S(P) = \bigcup_{p \in P}S_p\). On the cone \(\aaa^+\) we define the roots
\[
\alpha_{i,j}(a) = a_i - a_j,
\]
for \(i < j\) and notice that they are non-negative on \(\aaa^+\). Lyapunov and Falconer dimensions are special values of Lyapunov and Falconer functionals on \(\aaa^+\).

We define, for \(h \ge 0\) the Lyapunov functional on \(\aaa^+\) by
\begin{align*}L^P_h(a) &= \text{maximum of } \sum_{(i,j) \in S(P)}r_{i,j}
\\ &\hphantom{{}={}}\text{subject to }0 \le r_{i,j} \le 1
\quand \sum_{(i,j) \in S(P)} r_{i,j}\alpha_{i,j}(a) \le h.
\end{align*}
For \(r \ge 0\) we also define the Falconer functional on \(\aaa^+\) by
\begin{align*}F^P_r(a) &= \text{minimum of } \sum_{(i,j) \in S(P)}r_{i,j}\alpha_{i,j}(a)
\\ &\hphantom{{}={}}\text{subject to }0 \le r_{i,j} \le 1
\quand \sum_{(i,j) \in S(P)} r_{i,j} \ge r.
\end{align*}

\begin{lemma}[Duality between Falconer and Lyapunov functionals]\label{lyapunovvsfalconer}
Given \(a \in \aaa^+\), for all \(r,h \ge 0\) one has \(F_r^P(a) \le h\) if and only if \(L_h^P(a) \ge r\).

Furthermore, let \(d(a)\) be the number of pairs \((i,j) \) in \(S(P)\) with \(\alpha_{i,j}(a) = 0\).
Then, \(F^P_r(a) = 0 \) for \(r \in [0,d(a)]\) and \(r \mapsto F^P_r(a)\) is an increasing homeomorphism from \([d(a),D]\) to \([0,F_{D}^P(a)]\) where \(D = \# S(P)=\dim (\F_P)\), and \(h \mapsto L_h^P(a)\) is its inverse.
\end{lemma}
\begin{proof}
To establish the first claim observe that directly from the definitions, both conditions \(F_r^P(a) \le h\) and \(L_h^P(a) \ge r\) are equivalent to there being some choice of \(r_{i,j} \in [0,1]\) for \((i,j) \in S(P)\) such that
\[
\sum_{(i,j) \in S(P)} r_{i,j}\alpha_{i,j}(a) \le h\quand\sum_{(i,j) \in S(P)} r_{i,j} \ge r.
\]
Moreover, in computing \(F^P_r(a)\), we can always take \(r_{i,j} = 1 \) for the pairs \((i,j) \) in \(S(P)\) with \(\alpha_{i,j}(a) = 0\). It follows that \(F^P_r(a) = 0 \) for \(r \in [0,d(a)]\). Then, there exists \(\e >0\) with \(\alpha_{i,j}(a) > \e\) for all the other pairs \((i,j)\). It follows immediately that \(F^P_{r+s}(a) >F^P_r(a) + \e s\) for all \(d(a) \le r < r+s \le D\). Hence \(r \mapsto F^P_r(a)\) is increasing on \([d(a), D]\). Since \(F_{d(a)}^P(a) = 0\), \(r \mapsto F^P_r(a)\) it is an increasing homeomorphism as claimed.

Similarly, picking \(\e > 0\) small enough so \(\alpha_{i,j}(a) < \e^{-1}\) for all \((i,j) \in S(P)\), it~follows that \(L^P_{h + s}(a) > L^P_h(a) + \e s\) for all \(0 \le h < h+s < F_D^P(a)\). Hence, \(h \mapsto L_h^P(a)\) is an increasing homeomorphism as well.

Let \(f:I \to J\) and \(g:J \to I\) be the two homeomorphisms under consideration where \(I = [d(a),D]\) and \(J = [0,F_D^P(a)]\). By our first claim we have \(g(f(r)) \ge r\) for all \(r \in I\) and \(f(g(h)) \le h\) for all \(h \in J\). However, substituting \(r = f^{-1}(h)\), the first inequality implies \(g(h) \ge f^{-1}(h)\) for all \(h \in J\), from which we obtain \(f(g(h)) \ge h\) for all \(h \in J\). Hence, we~have \(f(g(h)) = h\) for all \(h \in J\) as required.
\end{proof}

Following \cite{DO80} and \cite{KY79}, we~define the Lyapunov dimension of \(\rho_*\mu\) on \(\F_P\) by
\begin{equation}\label{dimlyap:definition}\dim_{LY}^P(\rho_*\mu) \; := \; L_{h_\mu}^P (\lambda(\rho_*\mu)). \end{equation}
In the spirit of \cite{Fal88}, we~define the Falconer dimension \(\dim_F^P(\rho)\) of \(\rho\) relative to~\(\F_P\) as the critical parameter \(r \ge 0\) for convergence of the series
\begin{equation}\label{falconerseries} \Phi_\rho^P (r) \; := \; \sum_{\gamma \in \Gamma} \exp \bigl(-(F_r^P\circ\log S(\rho(\gamma)))\bigr).\end{equation}

\subsubsection{Proof of Theorem \ref{maintheorem}}

Theorem \ref{maintheorem} is a direct consequence of the three following results, all under the assumptions of Theorem \ref{maintheorem} and using the above notations:

\begin{theorem}[Upper bound on Lyapunov dimension]\label{theorema}
For all \(\mu \in \M\), one has
\[
\dim_{LY}^P(\rho_*\mu) \le \dim(\F_P) - \frac{M - 1}{2}.
\]
\end{theorem}

\begin{theorem}[Variational principle]\label{theoremb}
One has \begin{equation} \label{Lyapunov} \dim_F^P(\rho) = \sup_{\mu \in \M}\dim_{LY}^P(\rho_*\mu).\end{equation}
\end{theorem}

\begin{theorem}[Inequality between Minkowski and Falconer dimension]\label{theoremc}
One has
\[
\dim_B(\Lambda_\rho) \le \dim_F^P(\rho).
\]
\end{theorem}

The key observation behind Theorem \ref{theorema} is that the entropy \(h_\mu\) is realized as the Furstenberg entropy of the action of \(\rho (\G)\) on the Grassmannian \(\Gr(p,\R^d)\) for any \(p \in P\). This follows from the \(P\)-Anosov property and the identification of the Poisson boundary and geometric boundary for \(\mu\)-random walks on \(\Gamma\) (\cite{Kai00}). The inequality in Theorem \ref{theorema} is then a consequence of the inequality between Furstenberg entropy and the sum of the relevant exponents (see \cite{LL23} and Section \ref{prooftheorema}).

Concerning Theorem \ref{theoremb}, the fact that
\begin{equation}\label{LyapleqFalc} \dim_{LY}^P(\rho_*\mu) \; \leq \; \dim_F^P (G) \end{equation} is classical. We~recall the proof in Section \ref{proofLyap}.

The converse is due to Yuxiang Jiao, Jialun Li, Wenyu Pan and Disheng Xu (\cite{JLPX23}) in the case when the representation is Borel-Anosov. In \cite{JLPX23}, given \(\e>0\), the authors construct a free semi-group in \(\G\) that is rich enough that the uniform probability measure \(\mu_\e \) on the generators satisfies \(\dim_{LY}(\rho_*\mu_\e) \ge \dim_F(\rho) -\e\). The complete dominated splitting is used to control the almost additivity of the singular values.
Here, we~consider a general subset \(P\) and we assume that the representation \(\rho(\G)\) is Zariski dense in \(\SL(d,\R)\). In the spirit of \cite{gouezel}, \cite{random} and \cite{JLPX23}, we~prove in the appendix a general proposition about free sub-semigroups in hyperbolic groups (see the appendix for the definition of a quasi-geodesic set; see~also~\cite{yang} for a construction of free sub-semi-groups with convexity properties in a more general context.)
\begin{proposition}[$={}$Proposition \ref{semigrouplemma}]
Let \((\Gamma,\cS)\) be a non-elementary hyperbolic group with a finite symmetric generator \(\cS\). Then there exists a semigroup \(S \subset \Gamma\), a~finite set \(F \subset \Gamma\) and mappings \(L,R:\Gamma \to F\) with the following properties:
\begin{enumerate}
\item \(S\) generates \(\Gamma\) as a group,
\item \(S\) is quasi-geodesic, and
\item \(L(\gamma)\gamma R(\gamma) \in S\) for all \(\gamma \in \Gamma\).
\end{enumerate}
\end{proposition}
